\documentclass[aps,prb,onecolumn,nofootinbib,superscriptaddress]{revtex4-2}
\usepackage{amsmath,amssymb,amsthm,mathtools,bm}
\usepackage{booktabs}
\usepackage[colorlinks=true,linkcolor=blue,citecolor=blue,urlcolor=blue]{hyperref}
\usepackage{microtype}

\setcounter{tocdepth}{2}

\begin{document}

\title{Worker 06 Production Script and Metadata Audit}
\author{Worker 06}
\date{June 10, 2026}

\begin{abstract}
I independently inspected the production-facing script and notebook layer in
\texttt{scripts/}, \texttt{OSG/}, and current data-generation notebooks, excluding
paths containing \texttt{haining}, \texttt{haoyu}, or \texttt{erroneous}, and without
reading repo-synthesis worker, supervisor, or boss notes.  The strongest production
pattern is the recent GPU/OSG and Colab campaign style: call
\texttt{classA\_U1FGTN\_gpu.run\_markov\_circuit}, stream observables through
observers, avoid saving full covariance histories when unnecessary, and write
JSON/NPZ metadata with campaign manifests.  The main reliability risks are stale
legacy OSG scripts, incomplete canonical-entry metadata in one entropy campaign,
a likely cycle-label mismatch in one OSG merge script, and older CPU scripts that
produce scientific outputs without a separate manifest/index layer.
\end{abstract}

\maketitle
\tableofcontents

\section{Source Scope}

The inspection was read-only.  I did not edit files, execute notebooks, or load
large arrays.  I read repository policy first and used text inspection of source
files and notebook JSON only.

Included sources were:
\begin{itemize}
\item \texttt{scripts/}: production and exploratory CPU scripts for Markov
history generation, size sweeps, contour fitting, cache-derived generation, and
smoke tests.
\item \texttt{OSG/}: HTCondor submit files, OSG wrappers, parameter generators,
campaign runners, merge scripts, and bundled GPU helper/source copies.
\item Current data-generation notebooks under the Colab campaign trees, including
\texttt{colab\_small\_system\_testing/notebooks/data\_generation},
\texttt{colab\_charge\_fluctuations/notebooks/data\_generation},
\texttt{colab\_lyapunov/notebooks/data\_generation},
\texttt{colab\_regularized\_choi\_transfer\_matrix/notebooks/data\_generation},
the active \texttt{colab\_large\_entanglement\_scaling\_N20} generation notebooks,
and the partial-postselection GPU notebook.
\end{itemize}

Excluded by instruction were all paths containing \texttt{haining},
\texttt{haoyu}, or \texttt{erroneous}, and all
\texttt{repo\_synthesis/main\_pass} worker/supervisor/boss notes.

\section{Findings}

\subsection{Canonical Dynamics Usage}

The modern OSG runners use the canonical GPU entry point
\texttt{classA\_U1FGTN\_gpu.run\_markov\_circuit}.  The slope-vs-cycle and
slope-vs-system-size OSG scripts also store the literal key
\texttt{canonical\_dynamics\_entry\_point} in campaign configuration.  This is
well aligned with the repository policy.

The CPU scripts generally use \texttt{classA\_U1FGTN.run\_markov\_circuit}, which
is appropriate for local scripts.  However, some older generation scripts still
use older or adjacent APIs: \texttt{scripts/generate\_G\_from\_existing\_G0.py}
calls \texttt{run\_adaptive\_circuit}, while several older CPU scripts save raw
histories or figures without a manifest.  These are not necessarily wrong, but
they are weaker as production artifacts.

\subsection{Metadata and Campaign Traceability}

The best metadata pattern appears in newer GPU notebooks and OSG helpers: a
campaign-level configuration is serialized into JSON or \texttt{config\_json},
outputs are written atomically, and latest/campaign manifest pointers are
updated.  The notebook generation layer is especially strong in the small-system,
charge-fluctuation, Choi-transfer, and large-entanglement campaigns: it records
the canonical dynamics entry point, saved observables, covariance-history policy,
campaign IDs, run summaries, file sizes, or manifest pointers.

The OSG entropy-vs-time runner uses the canonical GPU method, but its saved
\texttt{config\_json} omits the explicit
\texttt{canonical\_dynamics\_entry\_point}.  This is a policy-level metadata gap:
the data are likely generated by the correct method, but downstream provenance
cannot prove that from the output alone.

\subsection{Campaign Generation}

The newer OSG campaign wrappers generate \texttt{params.json}, parse
\texttt{JOB\_COUNT}, export \texttt{N\_JOBS}, and submit with
\texttt{queue \$ENV(N\_JOBS)}.  This removes the old manual queue-count failure
mode.

The legacy root OSG path is much less reliable.  Its generator currently writes
four small-system jobs, but \texttt{OSG/batch\_gpu.submit} queues 100 jobs.  It
also contains a placeholder Docker image and saves only ndarray-valued return
objects, dropping explicit parameter and canonical-entry metadata.  This root OSG
path should be treated as stale unless manually repaired before use.

\subsection{Merge and Labeling Reliability}

There is a likely serious mismatch in the slope-vs-system-size OSG merge script.
The generator writes jobs with \texttt{cycles = Ny} and \texttt{fit\_cycle = Ny},
and the runner records those values in each block configuration.  The merge
script in \texttt{OSG/osg\_slope\_vs\_Ny} instead reconstructs
\texttt{cycles = Ny // 2} and \texttt{fit\_cycle = Ny // 2}.  Depending on helper
validation, this can either fail during checkpoint merge or silently produce
metadata labeling the wrong cycle.  This is the highest-impact issue found in the
inspected production layer.

The entropy-vs-time merge script checks cycle-array equality across blocks, which
is good, but it does not appear to validate that merged counts equal the intended
total sample count for every cycle and every \texttt{Ny}.  It reports the count at
the first cycle, so partial or uneven block recovery could be missed if not
noticed manually.

\subsection{Notebook Generation Layer}

Current Colab data-generation notebooks mostly follow the production conventions:
they import the GPU class rather than embedding full dynamics, check CUDA/A100
assumptions, call \texttt{run\_markov\_circuit}, serialize canonical-entry
metadata, and end with a Colab disconnect cell.  Several notebooks explicitly
state that covariance histories are not saved and instead stream observables,
which is the right storage model for large campaigns.

Some notebook JSON contains executed outputs.  That is not a runtime problem by
itself, but it increases diff size and can blur source/output provenance.  For
auditability, the source cells are still clear enough to identify the production
configuration and metadata scheme.

\section{Interpretation}

The repository has converged toward a sensible production architecture:
canonical dynamics are centralized in the CPU/GPU classes, large GPU campaigns
stream observables instead of saving unbounded histories, and new campaign
notebooks/scripts write enough metadata to reconstruct the scientific run.  The
remaining risk is not conceptual but operational.  Multiple campaign generations
coexist: modern manifest-based campaigns, intermediate OSG campaigns, and stale
root OSG scripts.  Without explicit deprecation or validation gates, a user could
submit an older path and obtain outputs that are hard to reproduce or even
mislabelled.

The single most important reliability implication is that downstream analysis
should prefer outputs containing an explicit canonical dynamics key, full
configuration JSON, sample/cycle counts, and manifest pointers.  Outputs from
older scripts should be treated as preliminary unless accompanied by external
run notes.

\section{Evidence Table}

\begin{table}[h]
\caption{Selected source evidence.  Line numbers are from text inspection only.}
\begin{ruledtabular}
\begin{tabular}{p{0.34\linewidth}p{0.18\linewidth}p{0.40\linewidth}}
Source & Lines & Evidence / implication \\
\hline
\texttt{REPO\_POLICY.md} & 1--26 & Requires production CPU/GPU simulations to use canonical \texttt{run\_markov\_circuit} entry points and record the entry point in metadata. \\
\texttt{OSG/scripts/run\_slope\_vs\_cycle\_block.py} & 65, 69--88, 183--197 & Defines \texttt{CANONICAL\_ENTRY\_POINT}, stores it in campaign config, calls GPU \texttt{run\_markov\_circuit} with \texttt{save=False} and a cycle observer. \\
\texttt{OSG/scripts/run\_slope\_vs\_cycle\_block.py} & 91--107, 201--214 & Validates block checkpoints and writes per-block status JSON. \\
\texttt{OSG/osg\_slope\_vs\_Ny/run\_slope\_vs\_system\_size.py} & 55, 72--83, 124--132 & Stores canonical-entry metadata and uses canonical GPU dynamics with streamed contour accumulation. \\
\texttt{OSG/osg\_slope\_vs\_Ny/generate\_params.py} & 41--57 & Generates full jobs with \texttt{cycles=Ny} and \texttt{fit\_cycle=Ny}. \\
\texttt{OSG/osg\_slope\_vs\_Ny/merge\_slope\_vs\_system\_size.py} & 54--58, 74--80 & Reconstructs \texttt{cycles=Ny//2} and \texttt{fit\_cycle=Ny//2}; this conflicts with the generator/runner convention. \\
\texttt{OSG/osg\_entropy\_vs\_time\_maxmix/run\_entropy\_vs\_time.py} & 98--114, 119--137 & Calls canonical GPU \texttt{run\_markov\_circuit}, but saved config lacks an explicit canonical-entry key. \\
\texttt{OSG/osg\_entropy\_vs\_time\_maxmix/merge\_entropy\_vs\_time.py} & 31--49 & Checks cycle-array equality across blocks and merges sums/sumsq/counts, but does not assert target count completeness for all cycles. \\
\texttt{OSG/osg\_slope\_vs\_Ny/submit\_campaign.sh} & 30--48 & Extracts \texttt{JOB\_COUNT}, exports \texttt{N\_JOBS}, and submits via environment-controlled queue length. \\
\texttt{OSG/osg\_slope\_vs\_Ny/all\_system\_sizes.submit} & 24--34 & Requests A100-class GPU resources and queues \texttt{\$ENV(N\_JOBS)}. \\
\texttt{OSG/run.sh} & 17--46 & Legacy root OSG runner reads params and calls GPU dynamics, but saves only ndarray results and no explicit run config. \\
\texttt{OSG/batch\_gpu.submit} & 1--2, 28 & Legacy submit file has a placeholder container image and hard-coded \texttt{queue 100}. \\
\texttt{OSG/generate\_params.py} & 48--58 & Legacy generator currently selects four jobs and instructs manual queue update. \\
\texttt{scripts/generate\_markov\_history\_samples.py} & 30--44, 99--113 & CPU allocation is handled through environment/affinity, and the canonical CPU Markov circuit is used with \texttt{save=True}. \\
\texttt{scripts/generate\_G\_from\_existing\_G0.py} & 55--77, 105--133 & Loads a cached snapshot and writes source sample/cycle metadata, but calls \texttt{run\_adaptive\_circuit} rather than the canonical Markov entry point. \\
\texttt{scripts/run\_markov\_size\_sweep\_integrated\_contour\_fit.py} & 17--31, 202--213 & CPU affinity/thread limits are set and canonical CPU \texttt{run\_markov\_circuit} is called. \\
\texttt{scripts/run\_markov\_size\_sweep\_integrated\_contour\_fit.py} & 267--308 & Saves a PDF figure but no manifest or machine-readable run summary. \\
\texttt{scripts/analyze\_covariance\_protocol\_histories\_characterization.py} & 325--340, 1139--1158 & Downstream analysis explicitly rejects unexpected canonical-entry metadata and writes an analysis manifest. \\
\texttt{colab\_large\_entanglement\_scaling\_N20/notebooks/run\_slope\_vs\_system\_size\_time\_scaling.ipynb} & 200--217, 398--406, 549--564, 587--588 & Notebook records canonical entry, uses GPU \texttt{run\_markov\_circuit}, writes a campaign manifest, and disconnects Colab runtime. \\
\texttt{colab\_small\_system\_testing/notebooks/data\_generation/run\_streaming\_covariance\_observables\_multi\_geometry.ipynb} & 239, 366--369, 440--450, 576--613, 641 & Records canonical entry, declares no covariance-history saving, streams observables, writes campaign metadata, and disconnects runtime. \\
\texttt{colab\_small\_system\_testing/notebooks/data\_generation/run\_purification\_maxmix\_entropy\_contours.ipynb} & 431--441, 467--518, 598--630, 646--647 & Uses GPU Markov dynamics with \texttt{save=False}, writes metadata/summary/manifest, and disconnects runtime. \\
\end{tabular}
\end{ruledtabular}
\end{table}

\section{Caveats}

This audit was source-level only.  I did not run scripts, submit OSG jobs, execute
notebooks, compare generated artifacts to manifests, or load large covariance
arrays.  I also did not inspect excluded paths or repo-synthesis notes.  Line
references in notebooks are JSON text-line references, not semantic notebook cell
numbers.  Some conclusions about stale scripts are based on internal inconsistency
and metadata absence, not on observed failed executions.

\section{Confidence}

Confidence is high for claims about source-level canonical-entry usage,
metadata keys, OSG queue generation, and the slope-vs-system-size cycle mismatch,
because these are directly visible in the inspected files.  Confidence is medium
for the production status of individual notebooks, since ``current'' was inferred
from active data-generation locations and campaign naming rather than from an
external registry.  Confidence is low for runtime behavior on OSG or Colab,
because no jobs or notebooks were executed.

\end{document}
